\naslov{Prvi integral enačbe}

Splošna rešitev enačbe $y' = f(x, y)$ je enoparametrična družina funkcij $y =
\phi(x, C)$.
Denimo, da obstaja taka funkcija $u(x, y) : [a, b] \times M \to \R$ na
razširjenem faznem prostoru, da zanjo velja $u(x) y(x, C) = \text{konst.}$ za
vsak $C$ (konstanta je lahko drugačna za različne $C$).
Taki funkciji pravimo \pojem{prvi integral enačbe}.
Recimo, da velja $\partial_y u \ne 0$.
Potem lahko iz enakosti izračunamo funkcijo $y(x, D)$, da velja $u(x, y(x, D)) =
D$.

\begin{trditev}
  Vsaka krivulja $y(x)$, ki je implicitno podana z enačbo $u(x, y) = D$, kjer je
  $u$ prvi integral enačbe $y' = f(x, y)$, je rešitev te enačbe.
\end{trditev}

\begin{proof}
  Naj bo $y_0(x)$ dana krivulja.
  Tedaj velja $u(x, y_0(x)) = D$.
  Če odvajamo po $x$, dobimo
  \[
	\partial_x u + y_0' \partial_y u = 0,
  \]
  torej
  \[
	y_0' = - \frac{\partial_x u}{\partial_y u}.
  \]
  Po drugi strani za vsako rešitev velja $y' = f(x, y)$,
  iz česar izpeljemo
  \[
	f(x, y) = -\frac{\partial_x u}{\partial_y u}.
  \]
  Sledi, da je $y_0$ res rešitev enačbe.
\end{proof}

Vsaka diferencialna enačba ima neskončno mnogo prvih integralov, vsakega lahko
še transformiramo s poljubno $\psi: \R \to \R$.

\vprasanje{Kaj je prvi integral enačbe? Kako iz njega dobiš rešitev enačbe?}

Imejmo dano vektorsko polje $F: \Omega \subseteq \R^2 \to \R^2$, podano z
\[
  F(x, y) =
  \begin{bmatrix}
	P(x, y) \\ Q(x, y)
  \end{bmatrix}
\]
Poiščimo družino krivulj, ortogonalnih na polje $F(x, y)$, in jih
parametrizirajmo z $\gamma(x) = (x, y(x))$.
Izpeljemo lahko pogoj
\[
  y'(x) = - \frac{P(x, y)}{Q(x, y)},
\]
kar je diferencialna enačba prvega reda.
Recimo, da je polje potencialno.
Tedaj obstaja taka funkcija $u: [a, b] \times \R \to \R$, da velja $\partial_x u
= P$ in $\partial_y u = Q$.
Krivulje, ki so ortogonalne na $\grad u$, so natanko izohipse ploskve $(x, y,
u(x, y))$. Za izohipso velja $u(x, y(x)) = C$, iz česar z odvajanjem izpeljemo
\[
  y' = - \frac{\partial_x u}{\partial_y u}.
\]
Potencial $u$ je torej prvi integral zgornje enačbe.
Ker lahko potencial poiščemo z integralom, enačbo v tem primeru znamo rešiti.

Če polje ni potencialno, imamo še vedno ortogonalne krivulje, torej še vedno
velja $y' = -\nicefrac{P}{Q}$.
Denimo, da je $y(x, C)$ splošna rešitev.
Če lahko poiščemo prvi integral $u$, mora obstajati funkcija $\lambda(x, y)$, za
katero je polje $(\lambda P, \lambda Q)$ potencialno, oziroma $\partial_x u =
\lambda P$ in $\partial_y u = \lambda Q$.
Taki funkciji pravimo \pojem{integrirujoči množitelj}.
Če ga lahko najdemo, lahko rešimo enačbo; tega pa v splošnem ne znamo.

\vprasanje{Kako poiščeš družino krivulj, pravokotnih na dano vektorsko polje
  $F$? Kaj je integrirujoči množitelj?}
